\documentclass[11pt]{article}
\usepackage{amsmath,amssymb}
\title{P1 N=6: full consolidated status after Rooms 1, 8, 11, 13--19 (6 reviewers)}
\date{2026-09-27}
\begin{document}
\maketitle

\section*{Corrected scope (Room 18/Reviewer BB)}
\textbf{Accumulation of poles at $\zeta_0=e(1/6)$ (and $\pm\sqrt{\cdot}$ for U,V) was NEVER open.}
\texttt{paper2a.tex} already proves this via closure from \texttt{thm:arc}/\texttt{cor:arcUV}
(its own text says so; the whole P1\_N6\_v2 series failed to cite it). The genuinely open question
is strictly stronger: \textbf{the direct, N=6-specific tie-ray statement} -- an explicit sequence
of zeros via an N=6 analogue of \texttt{prop:qi}'s own contour/Rouch\'e construction.

\section*{What is now PROVED, on the ray-lemma side of that direct statement}
Both halves of the separation-constant estimate that any such construction would need are secured,
each confirmed by independent adversarial review:
\begin{itemize}
\item \textbf{$\operatorname{Re}(x)\ge0$ (Room 17, confirmed by Reviewer BA after 2 prior genuine
errors caught in Rooms 14/16):} $d(x)=\min_{c^4=1}D(2\operatorname{Re}x,\cdot)\ge1-e^{-2\operatorname{Re}(x)}>0$,
band-uniform in $\theta$ near $\theta^\ast=0$, no seam.
\item \textbf{$\operatorname{Re}(x)\le0$ (Room 19, confirmed by Reviewer BC):} the paper's own
sign convention already routes around the danger -- $d'(x)=\min_{c^4=1}D(-2\operatorname{Re}x,\cdot)$
uses $a'=-2\operatorname{Re}(x)\ge0$ automatically on this domain, so the identical bound gives
$d'(x)\ge1-e^{2\operatorname{Re}(x)}>0$, by the same elementary argument, no new work needed.
\end{itemize}
Together: \textbf{the ray-lemma-based separation constant is positive everywhere on both parts of
the estimate, for the resonant class, uniformly in a band around $\theta^\ast=0$.} This closes the
specific obstruction (the $\theta^\ast=0$ degeneracy) that originally motivated this entire
investigation (\texttt{P1\_N6\_note.tex}, the very first room).

\section*{What is still genuinely missing}
\textbf{No N=6 contour construction exists anywhere in this repository} (confirmed independently
by Rooms 11, 13, 18, and Reviewer BC): no analogue of \texttt{prop:qi}'s $x_c$, $\Gamma_\pm^{(6)}$,
$\Pi_0^{(6)}$, or the Rouch\'e step that would turn the now-secured separation-constant bound into
an actual theorem. Two further, unattempted pieces sit behind that:
\begin{itemize}
\item The mod-6 Poisson-summation/triple-product bookkeeping producing $\Upsilon(x)$, $M'(x)$,
$K'(x)$ with correct numeric constants (\texttt{P1\_N6\_note.tex}'s original claim that this is
``mechanical'' has never been verified, though no obstruction to it has been found either).
\item The contour itself, and the Rouch\'e argument applied to it, which needs a genuine 2-D
neighborhood (not a ray) per Room 17's own finding about Rouch\'e's structural requirements.
\end{itemize}

\section*{Net assessment}
This 9-room, 6-reviewer sequence (Rooms 1, 8, 11, 13, 14, 16, 17, 18, 19) is the most thoroughly
adversarially-tested single thread of this entire multi-room effort. It caught 2 genuine errors in
sequence and 1 significant scope misunderstanding, ending with a real, confirmed, if narrow,
positive result (the separation constant, both sides) and a precisely isolated remaining task (the
contour construction + bookkeeping) for a future attempt. \textbf{Not integrated into paper2a}:
this does not close the direct tie-ray statement, only removes one obstruction on the path to it,
and the paper's own accumulation claim (which is what actually matters for the arc-density results
already published) was never affected either way.

\end{document}
